% Part: model-theory
% Chapter: interpolation
% Section: interpolation-proof

\documentclass[../../../include/open-logic-section]{subfiles}

\begin{document}

\olfileid{mod}{int}{prf}

\olsection{مبرهنة كريغ في الاستيفاء}

\begin{thm}[مبرهنة كريغ في الاستيفاء]
\ollabel{thm:interpol}
إذا كان $\Entails !A \lif !B$، وجدت !!a{sentence}~$!C$ بحيث
$\Entails !A \lif !C$ و$\Entails !C \lif !B$، وكان كل !!{constant}
و!!{function} و!!{predicate} (عدا~$\eq$) في $!C$ واقعًا في $!A$ و$!B$
معًا. وتسمى الـ!!{sentence}~$!C$ \emph{مستوفية} لـ$!A$ و$!B$.
\end{thm}

\begin{proof}
افترض أن $\Lang{L}_1$ لغة $!A$ وأن $\Lang{L}_2$ لغة $!B$. ولتكن
$\Lang{L}_0 = \Lang{L}_1 \cap \Lang{L}_2$. ولكل
$i \in \{0,1,2\}$، لتكن $\Lang{L}'_i$ ناتجة من $\Lang{L}_i$ بإضافة
عدد لا نهائي من الـ!!{constant}s الجديدة
$\Obj c_0, \Obj c_1, \Obj c_2, \dots$.

إذا كانت $!A$ غير قابلة للإرضاء، كانت
$\lexists[x][\eq/[x][x]]$ مستوفية. وإذا كانت $\lnot !B$ غير قابلة
للإرضاء (ومن ثم كانت $!B$ صالحة)، كانت
$\lexists[x][\eq[x][x]]$ مستوفية. ولذلك يجوز أن نفترض أيضًا أن $!A$
و$\lnot !B$ كلتيهما قابلتان للإرضاء.

لإثبات المقابل لمبرهنة الاستيفاء، افترض أنه لا توجد مستوفية لـ$!A$
و$!B$. وبعبارة أخرى، افترض أن $\{!A\}$ و$\{\lnot !B\}$ غير قابلتين
للفصل في $\Lang{L}_0$.

هدفنا توسيع الزوج $(\{!A\},\{\lnot!B\})$ إلى زوج غير قابل للفصل
أعظمي $(\Gamma^*,\Delta^*)$. ولتعدد $!A_0$ و$!A_1$ و$!A_2$ و\dots
الـ!!{sentence}s في $\Lang{L}'_1$، ولتعدد $!B_0$ و$!B_1$ و$!B_2$
و\dots الـ!!{sentence}s في~$\Lang{L}'_2$. نعرّف متتاليتين متزايدتين من
مجموعات الـ!!{sentence}s $(\Gamma_n,\Delta_n)$، حيث $n \ge 0$، كما
يأتي. ضع $\Gamma_0=\{!A\}$ و$\Delta_0=\{\lnot !B\}$. وبافتراض أن
$(\Gamma_n,\Delta_n)$ قد عرّف، عرّف $\Gamma_{n+1}$ و$\Delta_{n+1}$
كما يأتي:
\begin{enumerate}
\item إذا كانت $\Gamma_n \cup \{!A_n\}$ و$\Delta_n$ غير قابلتين
  للفصل في $\Lang{L}'_0$، فضع
  $\Gamma_{n+1}=\Gamma_n \cup \{!A_n\}$. وفوق ذلك، إذا كانت $!A_n$
  !!{formula} وجودية $\lexists[x][!S]$، فاختر !!{constant} جديدًا~$c$
  لا يقع في $\Gamma_n$ أو $\Delta_n$ أو $!A_n$ أو $!B_n$، وأضف
  $\Subst{!S}{c}{x}$ إلى $\Gamma_{n+1}$. وإذا كانت
  $\Gamma_n \cup \{!A_n\}$ و$\Delta_n$ قابلتين للفصل، فضع
  $\Gamma_{n+1}=\Gamma_n$.
\item إذا كانت $\Gamma_{n+1}$ و$\Delta_n \cup \{!B_n\}$ غير قابلتين
  للفصل في $\Lang{L}'_0$، فضع
  $\Delta_{n+1}=\Delta_n \cup \{!B_n\}$. وفوق ذلك، إذا كانت $!B_n$
  !!{formula} وجودية $\lexists[x][!S]$، فاختر !!{constant} جديدًا~$c$
  لا يقع في $\Gamma_{n+1}$ أو $\Delta_n$ أو $!A_n$ أو $!B_n$، وأضف
  $\Subst{!S}{c}{x}$ إلى $\Delta_{n+1}$. وإذا كانت $\Gamma_{n+1}$ و
  $\Delta_n \cup \{!B_n\}$ قابلتين للفصل، فضع
  $\Delta_{n+1}=\Delta_n$.
\end{enumerate}
وأخيرًا عرّف:
\begin{align*}
  \Gamma^* & = \bigcup_{n\ge 0} \Gamma_n, &
  \Delta^* & = \bigcup_{n\ge 0} \Delta_n.
\end{align*}
يمكننا الآن أن نثبت بالاستقراء المتزامن على $n$:
\begin{enumerate}
\item\ollabel{part-a} أن $\Gamma_n$ و$\Delta_n$ غير قابلتين للفصل في
  $\Lang{L}'_0$؛
\item\ollabel{part-b} أن $\Gamma_{n+1}$ و$\Delta_n$ غير قابلتين للفصل
  في $\Lang{L}'_0$.
\end{enumerate}
تعطينا \olref[sep]{lem:sep1} أساس \olref{part-a}. ولإثبات
\olref{part-b} نميز ثلاث حالات:
\begin{enumerate}
\item إذا كانت $\Gamma_0 \cup \{!A_0\}$ و$\Delta_0$ قابلتين للفصل،
  كان $\Gamma_1=\Gamma_0$، ولم تكن \olref{part-b} إلا
  \olref{part-a}؛
\item إذا كان $\Gamma_1=\Gamma_0 \cup \{!A_0\}$، كانت $\Gamma_1$
  و$\Delta_0$ غير قابلتين للفصل بحسب البناء.
\item تبقى الحالة التي تكون فيها $!A_0$ وجودية، بحيث
  $\Gamma_1=\Gamma_0 \cup \{\lexists[x][!S],\Subst{!S}{c}{x}\}$.
  وبحسب البناء، تكون $\Gamma_0 \cup \{\lexists[x][!S]\}$ و
  $\Delta_0$ غير قابلتين للفصل، ومن ثم، بحسب \olref[sep]{lem:sep2}،
  تكون $\Gamma_0 \cup \{\lexists[x][!S],\Subst{!S}{c}{x}\}$ و
  $\Delta_0$ غير قابلتين للفصل أيضًا.
\end{enumerate}
وهذا يتم أساس الاستقراء لـ\olref{part-a} و\olref{part-b} أعلاه. ننتقل
إلى خطوة الاستقراء. لأجل \olref{part-a}، إذا كان
$\Delta_{n+1}=\Delta_n \cup \{!B_n\}$، كانت $\Gamma_{n+1}$ و
$\Delta_{n+1}$ غير قابلتين للفصل بحسب البناء (حتى إذا كانت $!B_n$
وجودية، بفضل \olref[sep]{lem:sep2}). وإذا كان
$\Delta_{n+1}=\Delta_n$ (لأن $\Gamma_{n+1}$ و
$\Delta_n \cup \{!B_n\}$ قابلتان للفصل)، استعملنا فرض الاستقراء في
\olref{part-b}. ولأجل خطوة الاستقراء في \olref{part-b}، إذا كان
$\Gamma_{n+2}=\Gamma_{n+1} \cup \{!A_{n+1}\}$، كانت $\Gamma_{n+2}$
و$\Delta_{n+1}$ غير قابلتين للفصل بحسب البناء (حتى إذا كانت $!A_{n+1}$
وجودية، بفضل \olref[sep]{lem:sep2}). وإذا كان
$\Gamma_{n+2}=\Gamma_{n+1}$، استعملنا حالة الاستقراء لـ
\olref{part-a} التي أثبتناها للتو. وهذا يتم الاستقراء على
\olref{part-a} و\olref{part-b}.

يترتب أن $\Gamma^*$ و$\Delta^*$ غير قابلتين للفصل؛ وإلا لأعطى التراص
عددًا $n \ge 0$ تكون عنده $\Gamma_n$ و$\Delta_n$ قابلتين للفصل، خلافًا
لـ\olref{part-a}. وبوجه خاص، تكون $\Gamma^*$ و$\Delta^*$ متسقتين: إذ
لو كانت الأولى أو الثانية غير متسقة، لفصلت بينهما
$\lexists[x][\eq/[x][x]]$ أو $\lforall[x][\eq[x][x]]$، على الترتيب.

نبين الآن أن $\Gamma^*$ متسقة أعظميًا في $\Lang{L}'_1$، وكذلك
$\Delta^*$ في $\Lang{L}'_2$. ولأجل الأولى، افترض أن
$!A_n \notin \Gamma^*$ و$\lnot !A_n \notin \Gamma^*$ لبعض
$n \ge 0$. إذا كان $!A_n \notin \Gamma^*$، كانت
$\Gamma_n \cup \{!A_n\}$ قابلة للفصل عن $\Delta_n$، ومن ثم وجدت
$!C \in \Lang{L}'_0$ بحيث:
\begin{align*}
  \Gamma^* & \Entails !A_n \lif !C, &
  \Delta^* & \Entails \lnot !C.
\end{align*}
وبالمثل، إذا كان $\lnot !A_n \notin \Gamma^*$، وجدت
$!C' \in \Lang{L}'_0$ بحيث:
\begin{align*}
  \Gamma^* & \Entails \lnot !A_n \lif !C', &
  \Delta^* & \Entails \lnot !C'.
\end{align*}
وبالمنطق القضي، $\Gamma^* \Entails !C \lor !C'$ و
$\Delta^* \Entails \lnot (!C \lor !C')$، فتفصل $!C \lor !C'$
$\Gamma^*$ عن $\Delta^*$. وتثبت حجة مماثلة أن $\Delta^*$ أعظمية.

وأخيرًا نبين أن $\Gamma^* \cap \Delta^*$ متسقة أعظميًا في
$\Lang{L}'_0$. وهي متسقة بوضوح، لأنها تقاطع مجموعتين متسقتين. ولإثبات
الأعظمية، خذ $!S \in \Lang{L}'_0$. إن $\Gamma^*$ أعظمية في
$\Lang{L}'_1 \supseteq \Lang{L}'_0$، وبالمثل تكون $\Delta^*$ أعظمية
في $\Lang{L}'_2 \supseteq \Lang{L}'_0$. ومن ثم إما أن
$!S \in \Gamma^*$ أو $\lnot !S \in \Gamma^*$، وإما أن
$!S \in \Delta^*$ أو $\lnot !S \in \Delta^*$. فإذا كان
$!S \in \Gamma^*$ و$\lnot !S \in \Delta^*$، فصلت $!S$ بين
$\Gamma^*$ و$\Delta^*$. وإذا كان $\lnot !S \in \Gamma^*$ و
$!S \in \Delta^*$، فصلت $\lnot !S$ بينهما. ولذلك إما أن
$!S \in \Gamma^* \cap \Delta^*$ أو
$\lnot !S \in \Gamma^* \cap \Delta^*$، فيكون
$\Gamma^* \cap \Delta^*$ أعظميًا.

لما كانت $\Gamma^*$ متسقة أعظميًا، فلها نموذج $\Struct{M}'_1$ لا يضم
!!{domain}ه~$\Domain{M'_1}$ إلا العناصر $\Assign{c}{M'_1}$ التي تفسر
الـ!!{constant}s ويضمها جميعًا، كما في برهان مبرهنة الاكتمال
(\olref[fol][com][cth]{thm:completeness}). وبالمثل، لـ$\Delta^*$ نموذج
$\Struct{M}'_2$ يكون !!{domain}ه~$\Domain{M'_2}$ معطى بتفسيرات
$\Assign{c}{M'_2}$ للـ!!{constant}s.

لتكن $\Struct{M_1}$ ناتجة من $\Struct{M'_1}$ بحذف تفسيرات
الـ!!{constant}s والـ!!{function}s والـ!!{predicate}s في
$\Lang{L'_1} \setminus \Lang{L'_0}$، وبالمثل لـ$\Struct{M_2}$.
عندئذ يكون التطبيق $h \colon M_1 \to M_2$ المعرّف بـ
$h(\Assign{c}{M'_1})=\Assign{c}{M'_2}$ تشاكلًا في $\Lang{L}'_0$، لأن
$\Gamma^* \cap \Delta^*$ متسقة أعظميًا في $\Lang{L}'_0$ كما أثبتنا.
ويرجع ذلك إلى أن كل !!{sentence} من $\Lang{L}'_0$ إما أن تنتمي إلى
$\Gamma^*$ و$\Delta^*$ معًا، وإما ألا تنتمي إلى أي منهما: ولذلك
$\Assign{c}{M'_1} \in \Assign{P}{M'_1}$ إذا وفقط إذا كان
$\Atom{P}{c} \in \Gamma^*$، إذا وفقط إذا كان
$\Atom{P}{c} \in \Delta^*$، إذا وفقط إذا كان
$\Assign{c}{M'_2} \in \Assign{P}{M'_2}$. ويمكن إثبات سائر الشروط التي
تحققها التشاكلات بالطريقة نفسها.

نعرّف الآن نموذجًا~$\Struct{M}$ للـ!!{language}
$\Lang{L_1} \cup \Lang{L_2}$ كما يأتي:
\begin{enumerate}
\item الـ!!{domain}~$\Domain{M}$ هو $\Domain{M_2}$ نفسه، أي مجموعة كل
  العناصر $\Assign{c}{M'_2}$؛
\item إذا كان !!a{predicate}~$P$ في
  $\Lang{L_2} \setminus \Lang{L_1}$، فإن
  $\Assign{P}{M}=\Assign{P}{M'_2}$؛
\item إذا كان محمول~$P$ في $\Lang{L}_1\setminus \Lang{L}_2$، فإن
  $\Assign{P}{M}=h[\Assign{P}{M'_1}]$، أي إن
  $\tuple{\Assign{c_1}{M'_2},\dots,\Assign{c_n}{M'_2}} \in
  \Assign{P}{M}$ إذا وفقط إذا كان
  $\tuple{\Assign{c_1}{M'_1},\dots,\Assign{c_n}{M'_1}} \in
  \Assign{P}{M'_1}$.
\item إذا كان !!a{predicate}~$P$ في $\Lang{L}_0$، فإن
  $\Assign{P}{M}=\Assign{P}{M'_2}=h(\Assign{P}{M'_1})$.
\item تعالج !!^{function}s~$\Lang{L}_1 \cup \Lang{L}_2$، ومنها
  الـ!!{constant}s، بالطريقة نفسها.
\end{enumerate}

وأخيرًا يبيّن الاستقراء على الـ!!{formula}s أن $\Struct{M}$ توافق
$\Struct{M'_1}$ في جميع !!{formula}s~$\Lang{L_1}$، وتوافق
$\Struct{M'_2}$ في جميع !!{formula}s~$\Lang{L_2}$. وبوجه خاص،
$\Sat{M}{!A}$ و$\Sat{M}{\lnot!B}$، ومن ثم
$\not\Entails !A \lif !B$. وهذا يتم برهان مبرهنة كريغ في الاستيفاء.
\end{proof}

\end{document}
